\documentclass{amsart}
% For dealing with references we use the comment environment
\usepackage{verbatim}
\newenvironment{reference}{\comment}{\endcomment}
%\newenvironment{reference}{}{}
\newenvironment{slogan}{\comment}{\endcomment}
\newenvironment{history}{\comment}{\endcomment}

% For commutative diagrams we use Xy-pic
\usepackage[all]{xy}

% We use 2cell for 2-commutative diagrams.
\xyoption{2cell}
\UseAllTwocells

% We use multicol for the list of chapters between chapters
\usepackage{multicol}

% This is generally recommended for better output
\usepackage{lmodern}
\usepackage[T1]{fontenc}

% For cross-file-references

% Package for hypertext links:
\usepackage{hyperref}

% For any local file, say "hello.tex" you want to link to please

% Theorem environments.
%
\theoremstyle{plain}
\newtheorem{theorem}[subsection]{Theorem}
\newtheorem{proposition}[subsection]{Proposition}
\newtheorem{lemma}[subsection]{Lemma}

\theoremstyle{definition}
\newtheorem{definition}[subsection]{Definition}
\newtheorem{example}[subsection]{Example}
\newtheorem{exercise}[subsection]{Exercise}
\newtheorem{situation}[subsection]{Situation}

\theoremstyle{remark}
\newtheorem{remark}[subsection]{Remark}
\newtheorem{remarks}[subsection]{Remarks}

\numberwithin{equation}{subsection}

% Macros
%
\def\lim{\mathop{\mathrm{lim}}\nolimits}
\def\colim{\mathop{\mathrm{colim}}\nolimits}
\def\Spec{\mathop{\mathrm{Spec}}}
\def\Hom{\mathop{\mathrm{Hom}}\nolimits}
\def\Ext{\mathop{\mathrm{Ext}}\nolimits}
\def\SheafHom{\mathop{\mathcal{H}\!\mathit{om}}\nolimits}
\def\SheafExt{\mathop{\mathcal{E}\!\mathit{xt}}\nolimits}
\def\Sch{\mathit{Sch}}
\def\Mor{\mathop{\mathrm{Mor}}\nolimits}
\def\Ob{\mathop{\mathrm{Ob}}\nolimits}
\def\Sh{\mathop{\mathit{Sh}}\nolimits}
\def\NL{\mathop{N\!L}\nolimits}
\def\CH{\mathop{\mathrm{CH}}\nolimits}
\def\proetale{{pro\text{-}\acute{e}tale}}
\def\etale{{\acute{e}tale}}
\def\QCoh{\mathit{QCoh}}
\def\Ker{\mathop{\mathrm{Ker}}}
\def\Im{\mathop{\mathrm{Im}}}
\def\Coker{\mathop{\mathrm{Coker}}}
\def\Coim{\mathop{\mathrm{Coim}}}

% Boxtimes
%
\DeclareMathSymbol{\boxtimes}{\mathbin}{AMSa}{"02}

%
% Macros for moduli stacks/spaces
%
\def\QCohstack{\mathcal{QC}\!\mathit{oh}}
\def\Cohstack{\mathcal{C}\!\mathit{oh}}
\def\Spacesstack{\mathcal{S}\!\mathit{paces}}
\def\Quotfunctor{\mathrm{Quot}}
\def\Hilbfunctor{\mathrm{Hilb}}
\def\Curvesstack{\mathcal{C}\!\mathit{urves}}
\def\Polarizedstack{\mathcal{P}\!\mathit{olarized}}
\def\Complexesstack{\mathcal{C}\!\mathit{omplexes}}
% \Pic is the operator that assigns to X its picard group, usage \Pic(X)
% \Picardstack_{X/B} denotes the Picard stack of X over B
% \Picardfunctor_{X/B} denotes the Picard functor of X over B
\def\Pic{\mathop{\mathrm{Pic}}\nolimits}
\def\Picardstack{\mathcal{P}\!\mathit{ic}}
\def\Picardfunctor{\mathrm{Pic}}
\def\Deformationcategory{\mathcal{D}\!\mathit{ef}}

\setlength{\emergencystretch}{3em}
\begin{document}
\title[Stacks Project, Tag 09UF]{381 --- Finite-level pseudo-coherence descends under fppf coverings}
\maketitle
\noindent Copyright (C) 2005--2025 Johan de Jong. Stacks Project authors. Source: \href{https://stacks.math.columbia.edu/tag/09UF}{Tag 09UF}.

\noindent Editorial difficulty note (not author text). Difficulty H1: Descent of cohomology modules alone does not produce a finite-level complex resolution. The proof reduces the cover to one flat affine finite-presentation cover, lifts top-degree generators to a map from a shifted finite free sheaf, and inducts through its cone while retaining the precise pseudo-coherence bound..

\noindent Editorial provenance note (not author text). History: excerpt from revision \texttt{a0444\allowbreak 6e57e\allowbreak c1fbc\allowbreak 252a8\allowbreak 71afc\allowbreak ec775\allowbreak 2fb28\allowbreak 07b14}, perfect.tex, lines 2853--2898. Reference presentation was adapted; original-author context and core support blocks were retained or added. The mathematical wording is unchanged. Licensed under GNU FDL 1.2 or later; no invariant sections or cover texts. See ../licenses/COPYING. Context blocks reproduce author definitions and supporting results; they are not additional counted problems.

\begin{lemma}
\label{lemma-pseudo-coherent-descends-fppf}
Let $\{f_i : X_i \to X\}$ be an fppf covering of schemes. Let
$E \in D(\mathcal{O}_X)$. Let $m \in \mathbf{Z}$.
Then $E$ is $m$-pseudo-coherent if and only if each
$Lf_i^*E$ is $m$-pseudo-coherent.
\end{lemma}

\begin{proof}
Pullback always preserves $m$-pseudo-coherence, see
Cohomology, Lemma \href{https://stacks.math.columbia.edu/tag/09U7}{lemma pseudo coherent pullback (Tag 09U7)}.
Conversely, assume that $Lf_i^*E$ is $m$-pseudo-coherent for all $i$.
Let $U \subset X$ be an affine open. It suffices to prove that
$E|_U$ is $m$-pseudo-coherent. Since $\{f_i : X_i \to X\}$ is an
fppf covering, we can find finitely many affine open $V_j \subset X_{a(j)}$
such that $f_{a(j)}(V_j) \subset U$ and $U = \bigcup f_{a(j)}(V_j)$.
Set $V = \coprod V_i$.
Thus we may replace $X$ by $U$ and $\{f_i : X_i \to X\}$ by
$\{V \to U\}$ and assume that $X$ is affine and our covering
is given by a single surjective flat morphism $\{f : Y \to X\}$
of finite presentation.

\medskip\noindent
Since $f$ is flat the derived functor $Lf^*$ is just given by $f^*$ and $f^*$
is exact. Hence $H^i(Lf^*E) = f^*H^i(E)$. Since $Lf^*E$ is $m$-pseudo-coherent,
we see that $Lf^*E \in D^-(\mathcal{O}_Y)$. Since $f$ is surjective and flat,
we see that $E \in D^-(\mathcal{O}_X)$. Let $i \in \mathbf{Z}$ be the largest
integer such that $H^i(E)$ is nonzero. If $i < m$, then we are done. Otherwise,
$f^*H^i(E)$ is a finite type $\mathcal{O}_Y$-module by
Cohomology, Lemma \href{https://stacks.math.columbia.edu/tag/08DN}{lemma finite cohomology (Tag 08DN)}.
Then by Descent, Lemma \href{https://stacks.math.columbia.edu/tag/09UB}{lemma finite type descends fppf (Tag 09UB)}
the $\mathcal{O}_X$-module $H^i(E)$ is of finite type.
Thus, after replacing $X$ by the members of a finite affine open covering,
we may assume there exists a map
$$
\alpha : \mathcal{O}_X^{\oplus n}[-i] \longrightarrow E
$$
such that $H^i(\alpha)$ is a surjection. Let $C$ be the cone of $\alpha$
in $D(\mathcal{O}_X)$. Pulling back to $Y$ and using
Cohomology, Lemma \href{https://stacks.math.columbia.edu/tag/08CD}{lemma cone pseudo coherent (Tag 08CD)}
we find that $Lf^*C$ is $m$-pseudo-coherent. Moreover $H^j(C) = 0$
for $j \geq i$. Thus by induction on $i$ we see that $C$ is
$m$-pseudo-coherent. Using
Cohomology, Lemma \href{https://stacks.math.columbia.edu/tag/08CD}{lemma cone pseudo coherent (Tag 08CD)}
again we conclude.
\end{proof}
\end{document}
